\pdfoutput=1
\documentclass{article}

% ============================================================
% BaiZe-NeurIPS2027 — BaiZe MLLM 主线（从 BaiZe-ISEDA2027 拆分）
% 阶段：vision encoder pretrain -> MLLM alignment -> RL -> 评测
%       （通用12集 + 领域4集 + EDA Flow Bench 的 Multimodal 子集）
% TODO(内容重写)：0_abstract / 1_intro / 2_related_work / 3_architecture /
%                 7_conclusion 仍为 ISEDA 原文占位，需按 MLLM 主线重写；
%                 0_abstract 内 \begin{IEEEkeywords} 需改为 NeurIPS keywords 或删除。
% ============================================================

\usepackage[preprint]{neurips_2023}
\usepackage[utf8]{inputenc}
\usepackage[T1]{fontenc}
\usepackage{times}
\usepackage{url}
\usepackage{booktabs}
\usepackage{amsmath}
\usepackage{amssymb}
\usepackage{amsfonts}
\usepackage{bm}
\usepackage{graphicx}
\usepackage{array}
\usepackage{algorithm}
\usepackage{algpseudocode}
\usepackage[table]{xcolor}
\usepackage{xspace}
\usepackage{microtype}
\usepackage[hidelinks]{hyperref}

% IEEEkeywords stub（占位：抽象重写时移除）
\newenvironment{IEEEkeywords}{}{}

\newcommand{\llm}{BaiZe\xspace}

% TODO(标题定稿)：MLLM 主线标题占位
\title{BaiZe: From-Scratch Vision Encoder and Multimodal Grounding for Visual EDA Tasks}

% NeurIPS 匿名评审：作者信息留空
\author{}

\begin{document}

\maketitle

\begin{abstract}
Agentic assistants for electronic design automation (EDA) are already deployed in-house, but a correctness-critical slice of the flow---custom and analog design---resists them, because much of the correctness signal lives in \emph{visual} artifacts: layout, schematic, and waveform views that a text-only model cannot observe.
% 中文：面向电子设计自动化（EDA）的智能体助手已在公司内部署，但流程中与正确性紧密相关的一部分——定制与模拟设计——却难以用纯文本模型覆盖，因为大量正确性信号存在于版图、原理图与波形视图等视觉产物中，而这些是纯文本模型无法观测的。

We present the multimodal extension of \textbf{BaiZe}, an EDA agent built on the BaiZe language model~\cite{baiZe-llm}: a vision encoder pre-trained from scratch, aligned to the language backbone through a lightweight projector, and adapted to the semiconductor domain via reinforcement learning.
% 中文：我们提出 \textbf{BaiZe} 的多模态扩展——一个构建在 BaiZe 语言模型 \cite{baiZe-llm} 之上的 EDA 智能体：从零预训练一个视觉编码器，通过轻量投影层将其对齐到语言骨干，再通过强化学习适配到半导体领域。

A larger ViT-L/14 encoder reaches $68.2\%$ linear-probe on an internal protocol ($+3.9$\,pp over ViT-B/16), yet the aligned MLLM averages only $26.1\%$ over 11 multimodal benchmarks, indistinguishable from the smaller tower: encoder-quality gains do not transfer downstream at the $224$-px single-tile setting.
% 中文：更大的 ViT-L/14 编码器在内部协议上线性探测达 $68.2\%$（比 ViT-B/16 高 $+3.9$\,pp），但对齐后的 MLLM 在 11 个多模态基准上仅平均 $26.1\%$，与小塔不可分辨：编码器质量的增益在 224 像素单块设定下未迁移到下游。

We localise this gap and evaluate the full system on three tiers---12 general multimodal benchmarks, 4 domain (EDA) benchmarks, and the multimodal subset of an EDA flow benchmark---measuring whether a $196$-token-per-image budget suffices for detail-dense views, and whether domain reinforcement learning closes the transfer gap.
% 中文：我们定位这一落差，并在三个评测层上评估完整系统——12 个通用多模态基准、4 个领域（EDA）基准、以及 EDA 流程基准的多模态子集——度量每图 196 token 的预算是否足以应对细节密集的视图，以及领域强化学习能否弥合迁移落差。

We release the full recipe and an evaluation protocol, and report all measurements as \textbf{[TBD]}.
% 中文：我们公开完整配方与评估协议，所有测量结果以 \textbf{[TBD]} 表示（待填）。
\end{abstract}
\section{Introduction}

Agentic assistants for electronic design automation (EDA) are already deployed in-house, coordinating design flows through Flash-tier masters.
% 中文：面向电子设计自动化（EDA）的智能体助手已在公司内部署，通过 Flash 档 master 协调设计流程。
Yet a correctness-critical slice of the flow resists them: in custom and analog design, the signal that determines whether a layout ``works'' lives in \emph{visual} artifacts---layout, schematic, and waveform views---that a text-only model cannot observe.
% 中文：但流程中与正确性紧密相关的一部分却难以被它们覆盖：在定制与模拟设计中，决定版图是否「可用」的信号存在于版图、原理图与波形视图这类视觉产物中，而纯文本模型无法观测。
A subagent that can only read scripts and logs is blind to the very object it is designing, and must rely on an engineer to translate geometry into text---a bottleneck that scales with design volume.
% 中文：一个只能读取脚本与日志的子智能体，对其正在设计的对象本身是盲的，必须依赖工程师把几何信息翻译成文本——这一瓶颈随设计规模增长。

\paragraph{Contributions.}
We build the multimodal core of \textbf{BaiZe}---an EDA agent---on top of the BaiZe language backbone~\cite{baiZe-llm}, and ask a single question: \emph{can a from-scratch vision encoder make such a small model see the visual EDA artifacts that carry the correctness signal?}
% 中文：我们把 \textbf{BaiZe}——一个高性价比 EDA 智能体——的多模态内核构建在 BaiZe 语言骨干 \cite{baiZe-llm} 之上，并回答一个核心问题：\emph{一个从零训练的视觉编码器，能否让这样一个小模型「看见」承载正确性信号的视觉 EDA 产物？}

\begin{itemize}
\item \textbf{A from-scratch vision encoder.}
BaiZe pre-trains its own vision encoder rather than reusing an off-the-shelf one, and selects the tower from iso-parameter candidates at the $500$--$600$\,M scale.
% 中文：从零训练的视觉编码器。BaiZe 自研而非复用现成的视觉编码器，并在 $500$–$600$M 参数量级的同参数候选中选出编码器塔。
A ViT-L/14 tower reaches $68.2\%$ linear-probe on an internal protocol, $+3.9$\,pp above ViT-B/16, at a $196$-token-per-image budget.
% 中文：ViT-L/14 塔在内部协议上线性探测达 $68.2\%$，比 ViT-B/16 高 $+3.9$\,pp，且每图仅 196 token 预算。

\item \textbf{Token-bounded multimodal alignment.}
BaiZe connects the encoder to the frozen language backbone through a lightweight projector, following the staged curriculum of open multimodal assistants, and asks whether a $196$-token-per-image budget suffices for detail-dense layout and schematic views.
% 中文：成本受限的多模态对齐。BaiZe 通过轻量投影层把编码器连接到冻结的语言骨干，沿用开放多模态助手的分阶段课程，并检验每图 196 token 的预算是否足以应对细节密集的版图与原理图视图。
We show that encoder-quality gains do \emph{not} transfer downstream: the larger tower averages $26.1\%$ over 11 benchmarks, indistinguishable from the smaller one, and we localise the gap to the $224$-px single-tile setting and the alignment architecture.
% 中文：我们发现编码器质量的增益并未迁移到下游：更大的塔在 11 个基准上平均 $26.1\%$，与小塔不可分辨；我们把落差定位到 224 像素单块设定与对齐架构。

\item \textbf{Domain adaptation.}
A final reinforcement-learning stage adapts the multimodal model to proprietary semiconductor practice, learning team conventions that no public benchmark measures.
% 中文：领域适配。最后一个强化学习阶段把多模态模型适配到专有半导体实践，学习任何公开基准都无从衡量的团队约定。

\item \textbf{Three-tier evaluation.}
We evaluate end-to-end on 12 general multimodal benchmarks, 4 domain (EDA) benchmarks, and the multimodal subset of an EDA flow benchmark, under a single frozen execution-based protocol.
% 中文：三层评测。我们在 12 个通用多模态基准、4 个领域（EDA）基准，以及 EDA 流程基准的多模态子集上，按统一的冻结的、基于执行的协议进行端到端评测。
\end{itemize}
\section{Background and Related Work}

\subsection{Multimodal Vision-Language Models}
Vision-language models connect a vision encoder to a language model through a lightweight projector, enabling end-to-end visual understanding and instruction following.
% 中文：多模态视觉-语言模型。视觉-语言模型通过轻量投影层把视觉编码器连接到语言模型，实现端到端视觉理解与指令遵循。
LLaVA~\cite{liu2023llava} established visual instruction tuning---aligning a vision encoder to a frozen LLM via a linear projection followed by instruction fine-tuning---as a scalable recipe for open multimodal assistants.
% 中文：LLaVA \cite{liu2023llava} 把视觉指令微调——通过线性投影将视觉编码器对齐到冻结的 LLM，再进行指令微调——确立为开放多模态助手的可扩展配方。
LLaVA-OneVision~\cite{li2024llavaonevision} consolidates this line into a single open model spanning single-image, multi-image, and video tasks through a staged curriculum: stage-1 projector pre-training, stage-1.5 mid-training, and stage-2 instruction tuning.
% 中文：LLaVA-OneVision \cite{li2024llavaonevision} 将这一路线整合为单一开放模型，通过分阶段课程覆盖单图、多图与视频任务：stage-1 投影层预训练、stage-1.5 中期训练、stage-2 指令微调。
BaiZe follows this curriculum in stage~(iv) (Section~\ref{sec:mllm}), substituting a from-scratch vision encoder (stage~iii) for the frozen off-the-shelf encoder.
% 中文：BaiZe 在第 (iv) 阶段沿用这一课程（第 \ref{sec:mllm} 节），但把冻结的现成编码器替换为从零训练的视觉编码器（第 iii 阶段）。

\subsection{LLMs for Electronic Design Automation}
A growing line of work applies LLM agents to chip design tasks such as register-transfer-level generation, verification, and bug triage~\cite{liu2023chipnemo,blocklove2024chip}.
% 中文：面向电子设计自动化的 LLM。越来越多的工作把 LLM 智能体应用于芯片设计任务，如寄存器传输级生成、验证与缺陷分诊 \cite{liu2023chipnemo,blocklove2024chip}。
These applications motivate a cheap, always-on, agentic multimodal core that can run per-team or per-engineer; our work supplies the efficient multimodal backbone such systems require.
% 中文：这些应用呼唤一种廉价、常驻、可按团队或按工程师部署的智能体多模态内核；我们的工作正是提供此类系统所需的、高效的多模态骨干网络。

\section{Architecture and Pipeline Overview}\label{sec:overview}

BaiZe is a multimodal EDA agent assembled in three stages on top of the BaiZe language backbone~\cite{baiZe-llm}: (i) from-scratch pre-training of a custom vision encoder, (ii) alignment of the encoder and the backbone into a multimodal large language model, and (iii) semiconductor-domain adaptation via reinforcement learning.
% 中文：BaiZe 是一个构建在 BaiZe 语言骨干 \cite{baiZe-llm} 之上的、由三个阶段组装的多模态 EDA 智能体：(i) 从零预训练自研视觉编码器；(ii) 将编码器与骨干对齐为多模态大语言模型；(iii) 通过强化学习完成半导体领域适配。

This section fixes the shared scaffolding---the vision encoder, the alignment projector, and the global training environment---before each stage is detailed in its own section.
% 中文：本节先固定共享脚手架——视觉编码器、对齐投影层与全局训练环境——随后各阶段再在各自章节中详述。

The language backbone is the BaiZe $\approx$2.2B hybrid model, fixed and imported from its companion paper~\cite{baiZe-llm}; we do not re-train or re-describe it here, and keep it frozen throughout every stage.
% 中文：语言骨干是 BaiZe 约 22 亿参数的混合模型，从其姊妹论文 \cite{baiZe-llm} 固定导入；本文不重新训练也不赘述它，并在所有阶段中保持冻结。

\subsection{Vision Encoder}\label{sec:backbone-enc}
The vision encoder is pre-trained from scratch in stage~(i) and is selected from iso-parameter candidates at the $500$--$600$\,M scale; Section~\ref{sec:vision} reports the comparison and its outcome.
% 中文：视觉编码器。视觉编码器在第 (i) 阶段从零预训练，是在 $500$–$600$M 参数量级的同参数候选模型中选出的；第 \ref{sec:vision} 节报告该对比及其结果。

\subsection{Pipeline and Global Environment}\label{sec:pipeline}
Table~\ref{tab:pipeline} summarises the three-stage pipeline.
% 中文：流水线与全局环境。表 \ref{tab:pipeline} 汇总了三阶段流水线。
The vision encoder and alignment stages run on a single $8\times$H100 node (CUDA~12.8, PyTorch~2.8.0, Python~3.10); the aligned multimodal model uses a $196$-token-per-image budget throughout.
% 中文：视觉编码器与对齐阶段运行在单 $8\times$H100 节点上（CUDA 12.8、PyTorch 2.8.0、Python 3.10）；对齐后的多模态模型全程使用每图 196 token 的预算。

\begin{table}[t]
\centering
\caption{BaiZe multimodal pipeline overview.}
\label{tab:pipeline}
\begin{tabular}{llll}
\toprule
\textbf{Stage} & \textbf{Task} & \textbf{Status} & \textbf{Section} \\
\midrule
(i) & Vision encoder & complete & \ref{sec:vision} \\
(ii) & MLLM alignment & in progress & \ref{sec:mllm} \\
(iii) & Domain adapt.\ (RL) & design & \ref{sec:domain} \\
\bottomrule
\end{tabular}
\end{table}
\section{Vision Encoder Pretraining}\label{sec:vision}

Stage~(i) pre-trains a custom vision encoder from scratch; we summarise the outcome here and defer the design study to a companion report.
% 中文：第 (i) 阶段从零预训练一个自研视觉编码器；本节只汇报结果，设计研究另文详述。

\paragraph{Selected encoder.}
We adopt \textbf{OpenVision2}, a pure-attention ViT at width-$512$/depth-$30$ ($126.8$M parameters), trained from scratch with a bidirectional trunk at $224$px/$16$-px patches on a ${\approx}18.5$M-pair image--text corpus (CC12M $+$ Amshaker).
% 中文：选定编码器。我们采用 OpenVision2——宽度 $512$、深度 $30$ 的纯注意力 ViT，$126.8$M 参数，从零训练，双向主干，$224$ 像素/$16$ 像素块，在约 $18.5$M 图文对语料（CC12M + Amshaker）上训练。
The objective combines a global InfoNCE contrastive term with a caption-\emph{independent} masked-patch reconstruction term at a $1{:}1$ weight and a $0.6$ mask ratio.
% 中文：目标函数由全局 InfoNCE 对比项与一个与说明文字无关的掩码块重建项组成，权重 $1{:}1$，掩码比例 $0.6$。

\paragraph{Final result.}
On the frozen trunk with an ImageNet-1k linear probe, the encoder reaches $11.50\%\to\mathbf{19.76\%}$ top-1 from $5.1$M to $55.3$M samples (CC12M$+$Amshaker, $18.5$M unique pairs), and $20.27\%$ at $176$M samples on a $58.8$M-pair corpus (adding GPIC), with no saturation.
% 中文：最终结果。在冻结主干上用 ImageNet-1k 线性探测，编码器在 CC12M+Amshaker（$18.5$M 唯一对）上从 $5.1$M 到 $55.3$M 样本由 $11.50\%$ 升至 $\mathbf{19.76\%}$ top-1，在 $58.8$M 对语料（加入 GPIC）上于 $176$M 样本达到 $20.27\%$，均未见饱和。
This is $+12.4$\,pp above a contrastive-only baseline at the same sample budget, and it is the operating point used by stage~(ii).
% 中文：这比同样本量下仅用对比项的基线高 $+12.4$ pp，该点即第 (ii) 阶段使用的配置。
The $224/16$ resolution is retained because it is best on both loss and throughput; whether $196$ visual tokens suffice for detail-dense layout and schematic views is tested in stage~(ii) (Section~\ref{sec:mllm}).
% 中文：$224/16$ 分辨率被保留，因为它在损失与吞吐两方面都最优；$196$ 个视觉 token 是否足以应付细节密集的版图与原理图视图，将在第 (ii) 阶段检验（第 \ref{sec:mllm} 节）。

\paragraph{Encoder upgrade.}
A second tower---ViT-L/14 ($304$M) trained with open\_clip under the same objective but lower peak LR ($5{\times}10^{-4}$) and longer warm-up ($3000$ steps)---reaches $\mathbf{68.236\pm0.018\%}$ top-1 (5-NN $58.904\%$) on ImageNet-1k (Protocol~B, 3 seeds), versus $64.337\%$ / $55.956\%$ for the ViT-B/16 ($86$M) tower; the $+3.9$\,pp gap is a joint effect of tower size, LR, and warm-up, not a capacity-only attribution.
% 中文：编码器升级。第二个塔——ViT-L/14（304M），以 open_clip 在相同目标函数但更低峰值学习率（5×10⁻⁴）与更长预热（3000 步）下训练——在 ImageNet-1k 线性探测（Protocol B、3 seeds）上达到 68.236±0.018% top-1（5-NN 58.904%），对比 ViT-B/16（86M）塔的 64.337% / 55.956%；+3.9pp 的差距是塔规模、学习率与预热的联合效应，不能归因于容量单一变量。

\section{MLLM Alignment}\label{sec:mllm}

Stage~(ii) connects the vision encoder (Section~\ref{sec:vision}) to the language backbone via a lightweight projector, following the staged curriculum of open multimodal assistants; layout and schematic views account for $70.2\%$ of benchmark tasks~\cite{liu2026zhulong}, so a multimodal core shares the same execution-based protocol as the text-only stage~\cite{baiZe-llm}.
% 中文：第 (ii) 阶段通过一个轻量投影层把视觉编码器（第 \ref{sec:vision} 节）接到语言骨干网络上，沿用开放多模态助手的分阶段课程；版图与原理图视图占评测基准任务的 $70.2\%~\cite{liu2026zhulong}$，因此多模态内核可与纯文本阶段 \cite{baiZe-llm} 共享同一基于执行的协议。

\paragraph{Tasks and curriculum.}
Three task families each admit objective ground truth: (i) \emph{image$\to$script} (rendered view $\to$ PyAether/SKILL script, self-labelled by execution); (ii) \emph{script$\to$image} (predict geometry, self-labelled); (iii) \emph{layout VQA} (attribute questions, expert-labelled).
% 中文：任务与课程。三类任务各具客观标准答案：(i) \emph{图$\to$脚本}（渲染视图→PyAether/SKILL 脚本，执行自标注）；(ii) \emph{脚本$\to$图}（预测几何，执行自标注）；(iii) \emph{版图 VQA}（属性问题，专家标注）。
The first family shares the assertion sandbox of the text-only stage~\cite{baiZe-llm}, giving a common metric.
% 中文：第一类与纯文本阶段 \cite{baiZe-llm} 共享断言沙盒，提供统一度量。
We run stage-1 projector pre-training on captions, stage-1.5 mid-training on the three families, and stage-2 instruction tuning on the full mixture; ablations vary the projector type, tower freezing, and visual-token count.
% 中文：我们先用说明文字做 stage-1 投影层预训练，再用三类任务做 stage-1.5 中期训练，最后在全量混合上做 stage-2 指令微调；消融变量为投影层类型、塔冻结与否与视觉 token 数。
The $224/16$ setting ($196$ tokens/image) is retained; whether it transfers to detail-dense line drawings is what stage~(ii) tests.
% 中文：$224/16$（每图 196 token）设定被保留；其能否迁移到细节密集线稿正是第 (ii) 阶段所检验的。

\paragraph{Downstream results.}
Table~\ref{tab:mllm-bench} reports the Stage-2 step-3500 checkpoint on 12 lmms-eval benchmarks for ViT-L/14 ($304$M) and ViT-B/16 ($86$M).
% 中文：下游结果。表 \ref{tab:mllm-bench} 报告 Stage-2 step-3500 checkpoint 在 12 个 lmms-eval 基准上 ViT-L/14（304M）与 ViT-B/16（86M）的结果。
The 11 accuracy benchmarks average to $26.10\%$ vs.\ $26.12\%$ ($-0.02$\,pp, indistinguishable); despite the $+3.9$\,pp encoder advantage (Section~\ref{sec:vision}), downstream metrics are flat, suggesting the bottleneck lies not in tower capacity but in the $224$-px single-tile setting, the MLLM architecture, or weak proxy-to-downstream transfer.
% 中文：11 个准确率基准平均为 26.10% vs 26.12%（−0.02pp，不可分辨）；尽管编码器有 +3.9pp 优势（第 \ref{sec:vision} 节），下游指标持平，表明瓶颈不在塔容量，而在 224 像素单块设定、MLLM 架构或代理任务到下游的弱传递性。

\begin{table}[t]
\centering
\caption{MLLM downstream results (Stage-2 step~3500, \%; internal lmms-eval, single seed, $224$-px single-tile). MME$^\dagger$ is excluded from the 11-set mean.}
\label{tab:mllm-bench}
\setlength{\tabcolsep}{3pt}
\scriptsize
\begin{tabular}{@{}lcc|lcc@{}}
\toprule
Benchmark & ViT-L & ViT-B & Benchmark & ViT-L & ViT-B \\
\midrule
AI2D & 60.65 & 60.49 & CountBench & 10.59 & 9.98 \\
RealWorldQA & 45.36 & 46.41 & DocVQA & 7.35 & 7.44 \\
MMMU\_val & 44.00 & 42.67 & OCRBench & 1.60 & 1.60 \\
BLINK & 38.88 & 39.59 & InfoVQA & 16.18 & 16.59 \\
MMStar & 27.17 & 27.07 & ChartQA & 12.16 & 12.92 \\
MMMU\_Pro & 23.12 & 22.60 & \textbf{11-set avg} & \textbf{26.10} & \textbf{26.12} \\
\midrule
\multicolumn{6}{@{}l@{}}{MME$^\dagger$ (perc./cog./total): 513.75/204.29/718.04 \;\; 559.67/246.79/806.45} \\
\bottomrule
\end{tabular}
\end{table}

\section{Domain Adaptation: Waveform Understanding as a Probe Skill}\label{sec:domain}

Stage~(iii) adapts the multimodal core to semiconductor design practice. Rather than deferring this stage until a preference dataset matures, we treat \emph{waveform understanding}---reading simulation waveforms (VCD/FSDB) to answer questions about period, duty cycle, rise time, and signal shape---as a concrete, verifiable probe skill. It is the EDA sub-task where a multimodal model is most clearly distinguished from a text-only agent~\cite{baiZe-llm}: an engineer \emph{looks} at a waveform, and the relevant benchmarks down-grade precisely the visual judgements (eye-diagram aperture, edge margin, glitch detection) to text or Tcl when a vision channel is absent.
% 中文：第 (iii) 阶段把多模态内核适配到半导体设计实践。我们不再等待偏好数据集成熟，而是把「波形理解」——读取仿真波形（VCD/FSDB）来回答周期、占空比、上升时间与信号形状相关问题——当作一个具体的、可验证的探针技能。它是多模态模型与纯文本智能体 \cite{baiZe-llm} 区分最明显的 EDA 子任务：工程师「看」波形，而当视觉通道缺失时，相关基准恰好把眼图孔径、边沿余量、毛刺检测等视觉判断降级为文本或 Tcl。

\paragraph{Setup.}
We pilot the stage entirely on CPU, consistent with the deployment constraint that proprietary data stays on-premise on the BaiZe core~\cite{baiZe-llm}. The base model is Qwen3.5-2B (1.88B) with LoRA~\cite{hu2022lora} ($r{=}16$, ${\approx}6.4$M trainable parameters), trained in fp32 under \texttt{nice~-19} with a 16-thread cap, never touching the GPUs reserved for stage~(i)/(ii). A pure-Python waveform synthesizer generates sine/square/triangle/sawtooth signals with randomised amplitude, offset, frequency, phase and noise; the 16 uniform $(t,v)$ sample points are placed directly in the text prompt so that the language model can in effect ``see'' the shape. All numbers below are reproducible from the logged JSON in our experiment archive (train/eval use disjoint RNG seed ranges, so the cheat rate is zero).
% 中文：设置。我们完全在 CPU 上试点本阶段，与「专有数据须留在本地 BaiZe 内核 \cite{baiZe-llm}」的部署约束一致。基座为 Qwen3.5-2B（1.88B）+ LoRA \cite{hu2022lora}（r=16，约 6.4M 可训参数），fp32、nice -19、16 线程上限，绝不触碰预留给 (i)/(ii) 的 GPU。纯 Python 波形合成器生成正弦/方波/三角/锯齿信号，随机化幅度、偏置、频率、相位与噪声；16 个均匀 (t,v) 采样点直接放入文本 prompt，使语言模型事实上能「看见」形状。下列数字均可由实验存档中的 JSON 复现（训练/评测使用不相交 RNG seed 区间，故 cheat rate=0）。

\paragraph{Reward and task redesign.}
An advantage-variance diagnostic on the original objective revealed the bottleneck: with a 5-bin discrete reward and group size 8, 73.0\% of groups produced \emph{identical} rewards across all rollouts, so the GRPO~\cite{shao2024deepseekmath} advantage was zero and the gradient carried no signal; the zero-advantage share was \emph{not} the ``all-zero'' failure mode (0.0\%) but the ``all-same-bin'' mode (15.0\%). Per-target, \emph{period} was the only target with usable signal (std$>$0 in 78.3\% of groups), whereas \emph{rise\_time} was 100\% zero-variance and \emph{duty\_cycle} 91.7\%---a capability gap, not a compute gap. We therefore redesigned on two axes: (i) a \emph{continuous} reward $R=\exp(-\mathrm{rel\_err}/\tau)$ replacing the discrete bins, and (ii) \emph{task redefinition} from exact-number regression to a discriminative intermediate task (waveform \emph{type classification}, then \emph{frequency comparison}), which is decidable, informative for downstream use, and far more learnable for a 2B adapter. The adaptation methods are rejection-sampling fine-tuning (RFT, in the STaR/ReST family~\cite{zelikman2022star,gulcehre2023rest}) and outcome-only GRPO with a KL-free group baseline.
% 中文：奖励与任务重设计。对原目标做 advantage 方差诊断揭示了瓶颈：5 档离散奖励 + group size 8 时，73.0\% 的组内所有 rollout 获得相同奖励，故 GRPO \cite{shao2024deepseekmath} 的 advantage=0、梯度无信号；零 advantage 并非「全 0」失效模式（0.0\%），而是「全同档」模式（15.0\%）。按目标拆分，period 是唯一有可用信号的（78.3% 组 std>0），而 rise\_time 100% 零方差、duty\_cycle 91.7%——这是能力缺口而非算力缺口。故我们在两条轴上重设计：(i) 连续奖励 $R=\exp(-\mathrm{rel\_err}/\tau)$ 取代离散分档；(ii) 任务重定义，由精确数字回归改为可判定的中间任务（波形类型分类，继而频率比较）。适配方法为拒绝采样微调（RFT，属 STaR/ReST 家族 \cite{zelikman2022star,gulcehre2023rest}）与 outcome-only、无 KL 的 group-baseline GRPO。

\paragraph{Protocol and falsification.}
Every arm is pre-registered before measurement: four conditions are evaluated on a fixed held-out set with $\geq 3$ random seeds. \textbf{A}~=~fresh LoRA baseline; \textbf{B}~=~A after SFT; \textbf{C}~=~B after the adaptation step (RFT or GRPO); \textbf{D}~=~control (random-filter RFT or random-reward GRPO). Success requires $C{>}B$, $C{>}D$, a McNemar test on $C$ vs $B$ with $p{<}0.05$, and a bootstrap 95\% CI on $C{-}B$ that excludes zero. Negative results are reported as-is; a CPU pilot is never extrapolated to a GPU claim.
% 中文：协议与证伪。每个臂在测量前预注册：在固定留出集上以 $\geq 3$ 个随机种子评测四条件。A=fresh LoRA 基线；B=SFT 后的 A；C=适配步骤后的 B；D=控制组。成功需 $C{>}B$、$C{>}D$、McNemar $p{<}0.05$、且 bootstrap 95\% CI 排除 0。负结果如实记录；CPU 试点从不外推为 GPU 结论。

\paragraph{Results.}
Table~\ref{tab:domain-rl} summarises seven consecutive design iterations. The single statistically significant positive is the type-classification task with RFT (\emph{v5}): accuracy climbs $A{=}16.0\%\to B{=}37.0\%\to C{=}53.0\%$, with McNemar $C$-vs-$B$ $p{=}0.0009$ and bootstrap CI $[+7\%,+25\%]$, while the random-filter control $D{=}32.0\%$ falls \emph{below} $B$---the gain is driven by the correctness filter, not by data volume. The improvement concentrates on the visually distinctive \emph{square} class (84\%); \emph{triangle}/\emph{sawtooth} remain near zero, a persistent capability gap. Switching the adapter from RFT to GRPO on the same task (\emph{v6}) drops to 40\% (McNemar $p{=}0.267$, NS); adding a process-reward term (\emph{v7}) reaches $49.3\%\pm6.4\%$ and passes 6/6 pre-registered criteria yet fails the McNemar test because the seed-0 baseline runs high. Redefining the task to \emph{frequency comparison} (\emph{v8}) yields a mixed $52.5\%\pm6.1\%$ (significant for one seed only), and three further iterations (\emph{v9}--\emph{v11})---including a generative-augmentation route---produce no advantage over the direct arm (\emph{v11}: direct 22.2\% vs.\ augment 20.8\% vs.\ synthesis 13.2\%, $0/6$ criteria). We read this honestly: \emph{one} task--method pairing gives a significant, controlled gain, and it does \emph{not} transfer to the harder regression or frequency tasks within the CPU budget.
% 中文：结果。表 \ref{tab:domain-rl} 汇总连续七轮设计迭代。唯一统计显著正结果是 RFT 上的类型分类（v5）：$A=16.0\%\to B=37.0\%\to C=53.0\%$，McNemar $p=0.0009$、CI $[+7\%,+25\%]$，随机过滤控制 $D=32.0\%$ 低于 $B$——增益来自正确性过滤而非数据量。提升集中在方波类（84\%）；三角/锯齿仍接近零。同任务换 GRPO（v6）降至 40%（NS）；加过程奖励（v7）达 $49.3\%\pm6.4\%$，6/6 判据通过但 McNemar 不显著。频率比较（v8）混合 $52.5\%\pm6.1\%$（仅一 seed 显著），其后 v9–v11（含生成增强）对 direct 无优势（v11：22.2/20.8/13.2，0/6 判据）。如实解读：仅一个「任务×方法」组合给出受控显著增益，且在 CPU 预算内不迁移到更难任务。

\begin{table}[t]
\centering
\caption{Stage~(iii) CPU pilot on Qwen3.5-2B + LoRA. Accuracy (\%) per pre-registered arm; NS = McNemar $p{\geq}0.05$. All runs: train/eval disjoint seeds $\Rightarrow$ cheat rate $=0$. ``--'' = not applicable to that task.}
\label{tab:domain-rl}
\setlength{\tabcolsep}{4pt}
\scriptsize
\begin{tabular}{@{}llcccccl@{}}
\toprule
Iter & Task (method) & $A$ & $B$ & $C$ & $D$ (ctrl) & McNemar $C{ vs.}B$ & Verdict \\
\midrule
v5 & Type cls., 4-way (SFT$\to$RFT) & 16.0 & 37.0 & \textbf{53.0} & 32.0 & $p{=}0.0009$ & \textbf{Significant} \\
v6 & Type cls. (SFT$\to$GRPO) & -- & 35.0 & 40.0 & -- & $p{=}0.267$ & NS \\
v7 & Type cls. (GRPO+process) & -- & ${\sim}40$ & $49.3{\pm}6.4$ & -- & NS & 6/6 crit., NS test \\
v8 & Freq.\ compare (SFT$\to$GRPO) & 0.0 & 45.0 & $52.5{\pm}6.1$ & 45.0 & $p{=}0.031$ (s0) & Mixed \\
v9 & Freq.\ compare (GRPO) & 0.0 & 49.5 & 49.5 & -- & $p{=}1.0$ & No change \\
v10 & Freq.\ compare (GRPO) & 46.2 & 54.9 & $51.6{\pm}3.1$ & 53.8 & $p{=}1.0$ & NS \\
v11 & Regen.\ vs.\ direct & 22.2 & -- & 20.8 / 13.2$^*$ & -- & -- & $0/6$, negative \\
\bottomrule
\end{tabular}
\\[2pt]
\footnotesize $^*$v11 $C$ columns are \emph{direct / augment / synthesis} arms; generative augmentation is worse than direct.
\end{table}

\paragraph{Where vision would actually help.}
To bound the value of a true multimodal (image) channel, we audited the waveview-related cases of an EDA-flow benchmark~\cite{liu2026zhulong}: of $37$ waveform-related cases ($30$ failed), $4$ ($\alpha{=}13.3\%$, optimistic $\alpha{+}\beta{\approx}36\%$) fail \emph{directly because} the model cannot see the waveform shape---eye-diagram aperture, rise-edge margin, root-cause diff, functional inference from a plot. This audit was conducted under a benchmark version that \emph{downgraded} the visual judgement of $7$ evaluated cases to text/Tcl (all passing); downgrading the very capability under test changes the ruler, so that version has since been corrected---need-to-see-image cases are retained in original form and a system without multimodal ability is now \emph{expected} to fail them. The $\alpha$/$\beta$ cases therefore become the measurement substrate for multimodal gain, defined as the \emph{baseline-fail~$\to$~with-vision-pass} delta, which the cross-line re-evaluation below targets. The implication for stage~(iii) is that the image-channel ceiling must be re-measured under the corrected benchmark rather than read off the downgraded numbers, and that task/reward design (above) together with a waveform-VQA set built from the $\alpha$/$\beta$ cases are the high-ROI levers.
% 中文：视觉真正能帮上忙之处。为界定真实多模态（图像）通道的价值，我们审计了某 EDA-flow 基准 \cite{liu2026zhulong} 的 waveview 相关用例：37 个波形相关用例（30 个 failed）中，4 个（$\alpha=13.3\%$，乐观 $\alpha+\beta\approx36\%$）是「直接因为模型看不到波形形状」而失败——眼图孔径、上升沿余量、根因 diff、由图作功能推断。该审计是在「把 7 个已评测用例的视觉判断降级为文本/Tcl（全通过）」的基准版本下做的；降级被测能力等于换了尺子，故该版本此后已被纠正——需看图的用例恢复原形态，无多模态能力的系统现在「应当」在这些题上 fail。于是 $\alpha$/$\beta$ 用例成为多模态增益的度量基底，增益定义为「基线 fail→上视觉后 pass」的差值，正是下方跨线再评测的目标。对第 (iii) 阶段的含义是：图像通道天花板须在纠正后的基准上重测，而非取自降级数字；且任务/奖励设计（上文）+ 用 $\alpha$/$\beta$ 用例构造 waveform-VQA 评测集才是高 ROI 杠杆。

\paragraph{Honest status.}
Within a CPU-only budget, stage~(iii) delivers one controlled, significant result (type classification, $+16$\,pp over SFT, $p{=}0.0009$, cheat rate $0$) and a sequence of non-significant or negative results that we report rather than suppress; the binding constraint is reward/task definition, not compute, so waiting for GPUs would not have changed the qualitative outcome. Two items remain open and are stated as such rather than estimated: (i) scaling the significant pairing to GPU-scale GRPO, and (ii) the cross-line re-evaluation that wraps Qwen3.5-2B as a VQA tool and re-runs the $\alpha$/$\beta$ waveview cases against a text-only baseline---the dependency chain whose first link is the audit above.
% 中文：诚实状态。在仅 CPU 预算内，第 (iii) 阶段交付一个受控显著结果（类型分类，较 SFT +16pp，$p=0.0009$，cheat rate 0）以及一连串如实报告的非显著或负结果；约束瓶颈是奖励/任务定义而非算力，故等待 GPU 不会改变定性结论。两项仍开放，据实陈述而非估算：(i) 把显著组合扩展到 GPU 规模 GRPO；(ii) 跨线再评测——把 Qwen3.5-2B 封装为 VQA 工具、对 $\alpha$/$\beta$ waveview 用例相对纯文本基线重跑——其依赖链第一环即上文的审计。

\section{Conclusion}

We built the multimodal core of BaiZe on the frozen BaiZe language backbone~\cite{baiZe-llm}: a vision encoder pre-trained from scratch, aligned to the backbone through a lightweight projector, and adapted to the semiconductor domain via preference-based reinforcement learning.
% 中文：我们在冻结的 BaiZe 语言骨干 \cite{baiZe-llm} 之上构建了 BaiZe 的多模态内核：从零预训练的视觉编码器，经轻量投影层对齐到骨干，再通过基于偏好的强化学习适配到半导体领域。

Through a controlled encoder comparison, we found that a larger ViT-L/14 tower ($68.2\%$ linear-probe, $+3.9$\,pp over ViT-B/16) does \emph{not} translate its advantage downstream: the aligned MLLM averages $26.1\%$ over 11 benchmarks, indistinguishable from the smaller tower. We attribute the gap to the $224$-px single-tile setting and the alignment architecture rather than to tower capacity.
% 中文：通过受控的编码器对比，我们发现更大的 ViT-L/14 塔（线性探测 $68.2\%$，比 ViT-B/16 高 $+3.9$\,pp）并未把优势迁移到下游：对齐后的 MLLM 在 11 个基准上平均 $26.1\%$，与小塔不可分辨。我们把这一落差归因于 224 像素单块设定与对齐架构，而非塔容量。

Stages~(ii) and~(iii) are evaluated across 12 general multimodal benchmarks, 4 domain benchmarks, and the multimodal subset of an EDA flow benchmark. Stage~(iii) is reported here as a CPU-only pilot (Section~\ref{sec:domain}): it yields one controlled, statistically significant gain on a waveform type-classification probe ($+16$\,pp over SFT, $p{=}0.0009$, cheat rate $0$) alongside a sequence of honestly reported non-significant or negative iterations, and an audit that bounds the ceiling of an image channel on the current benchmark. Scaling the significant pairing to GPU-scale RL, and the cross-line VQA-tool re-evaluation of the waveview cases, remain open and are stated as such rather than estimated.
% 中文：第 (ii) 与 (iii) 阶段在 12 个通用多模态基准、4 个领域基准与 EDA 流程基准的多模态子集上评测。第 (iii) 阶段在此作为仅 CPU 试点报告（第 \ref{sec:domain} 节）：它在波形类型分类探针上给出一个受控统计显著增益（较 SFT +16pp，$p=0.0009$，cheat rate 0），同时如实报告一连串非显著或负的迭代，以及一份界定「图像通道在当前基准上天花板」的审计。把显著组合扩展到 GPU 规模 RL、以及对 waveview 用例的跨线 VQA 工具再评测，仍属开放，据实陈述而非估算。

The immediate next step is closing the encoder-to-downstream transfer gap---through higher-resolution multi-tile input and domain-specific alignment---toward a deployable multimodal core for visual EDA tasks.
% 中文：下一步工作是弥合「编码器到下游」的迁移落差——通过更高分辨率的多块输入与领域专项对齐——目标是得到一个可用于视觉 EDA 任务的可部署多模态内核。


\bibliographystyle{plain}
\bibliography{reference}

\end{document}